\subsection{The output of the application, panel by panel}
The sprinkler layer is organised in five parts: the site, the sprinkler and its layout, the way the system is run, the farm and its pipes, and the results with the course checks. The case shown is a solid-set system of impact sprinklers with 4.4 mm nozzles at 300 kPa, spaced 12 m along laterals 18 m apart, on 5 ha of coarse loam with a wind of 3 m\,s$^{-1}$; it applies 72 mm per irrigation.

The first part sets the conditions that limit the sprinkler. The soil tiles give the maximum application rate of each texture at the slope chosen (25 mm\,h$^{-1}$ for coarse loam on a 1.0 per cent slope), which is the rate the sprinkler must not exceed. The wind speed (3.0 m\,s$^{-1}$) enters the drift loss, and the surface storage (5.0 mm) is the depth of water the soil surface can hold before it runs.

\appfig{spr_site}{0.96}{Part A, your site: crop, mulch, soil type with the largest application rate it accepts for the slope, ground slope, wind and surface storage.}

The buttons choose the system (solid-set, hand-move, centre pivot, linear move, travelling gun) and the sprinkler (impact, rotating spray, micro-sprinkler, big gun) and the square or triangular layout. The numerical boxes set the nozzle diameter (4.4 mm), the trajectory angle (25$^\circ$), the pressure (300 kPa) and the spacing along and between laterals. The live view on the left draws the wetted circle (throw radius 15.7 m) and the one on the right the head of the sprinkler with the jet and the droplets; the text under it gives the flow (21.2 L\,min$^{-1}$), the jet speed (14 m\,s$^{-1}$), the biggest drop (3.8 mm) and the loss to wind and air (3.8 per cent).

\appfig{spr_layout}{0.96}{Part B, the sprinkler and its layout: system type, sprinkler type, layout, nozzle, trajectory angle, pressure and spacing, with the live view of the throw.}

Part C sets the depth per irrigation (72 mm, from the DSS) and the operating hours per day (16 h); the run time of 805 min is computed. Part D collects the area (5.00 ha), the lateral and mainline lengths, the farm supply (40 L\,s$^{-1}$) and the lift (12 m) from which the pipes and the pump are sized.

\appfig{spr_run}{0.96}{Parts C and D, how you run it and the farm with its water supply and pipes.}

The plan shows the sprinklers on the grid and the circles of the spray patterns overlapping. The colours give the depth infiltrated at each point: green is the depth needed, red-brown is short, blue is more than needed and brown hatching is water running off. The overlap of the patterns is the origin of the Christiansen uniformity of this chapter

\appfig{spr_plan}{0.96}{Layout in plan: depth infiltrated so far, from the overlapping spray patterns of the sprinklers.}

The section along the lateral shows the water that arrives at each point (the applied depth, including the drift loss), the depth that infiltrates and the volume of the soil that receives it; below the real texture of the soil the coloured column is the water that infiltrated, the depth divided by the water the soil can hold. The profile of the day at the bottom gives the trigger and field capacity for reference.

\appfig{spr_lateral}{0.96}{Along the lateral: depth applied, depth infiltrated and where it ends up in the soil, with the soil water of the day.}

The curve is the infiltration rate of the soil with time (the derivative of the Kostiakov intake of this chapter) and the horizontal line is the application rate of the sprinklers (5.9 mm\,h$^{-1}$, or 5.7 mm\,h$^{-1}$ after the drift loss). The soil intake at 30 min (19.2 mm\,h$^{-1}$) is well above the application rate, so no runoff occurs; the table lists ponding, runoff, surface storage and pattern uniformity.

\appfig{spr_intake}{0.96}{Soil intake against the application rate: when the soil cannot keep up, water runs off.}

The cards and the bar split the water that leaves the sprinklers into the part stored in the roots (91 per cent), the deep percolation, the drift and evaporation and any runoff; they are the quantities $E_a$, $E_s$ and $CU$ of the chapter.

\appfig{spr_result}{0.96}{Result for this design: application efficiency 90.6 per cent, root zone refilled 99.8 per cent, Christiansen uniformity 96.2 per cent, runoff 0, deep percolation 5.6 per cent, drift and evaporation 3.8 per cent.}

The layer applies the rules of the course as ten checks: the spacing along the lateral and between laterals against the wetted radius, the wind limit, the application rate against the maximum for the soil and against the saturated conductivity, the Christiansen uniformity, the lateral head loss (at most 20 per cent of the working head), the mainline gradient, the time to irrigate the whole field within the interval and the flow within the farm supply. A green tick means the rule is met, with the actual value written beside it.

\appfig{spr_checks}{0.96}{Does your design pass? The course checks, ten of ten satisfied.}

The teaching part of the layer begins with the kinds of system (solid-set, hand-move, centre pivot, linear move and travelling gun), each with its plan and its range of application efficiency (70 to 75 per cent for solid-set, 65 to 70 per cent for the self-propelled types), and shows where this design falls on the scale.

\appfig{spr_kinds}{0.96}{How a sprinkler works: the kinds of sprinkler system and their efficiency ranges.}

The panel names the parts of an impact sprinkler (the swing arm, the spring, the nozzle, the bearing and the riser) and classifies sprinklers by their operating pressure: low pressure (6 to 20 m), medium pressure (20 to 45 m, which includes this design) and high pressure (45 to 100 m, the guns). It repeats the discharge law $q=C_d\,a\sqrt{2gh}$ with $q=k\,h^{0.5}$ and the throw of the jet.

\appfig{spr_parts}{0.96}{The sprinkler: parts and pressure class, with the head of a rotary impact sprinkler.}

The farthest throw is at about 28 to 30 degrees; low angles (8 degrees) go under foliage and drift less. The panel shows how small droplets fall near the sprinkler and large ones far away (0.5 to 4 mm).

\appfig{spr_jet}{0.96}{The jet: trajectory angle, throw radius and droplet size.}

